\documentclass[10pt,a4paper]{amsart}
\usepackage[margin=19mm]{geometry}
\usepackage{amsmath,amssymb,mathtools}
\usepackage[scr=rsfs,cal=euler]{mathalfa}
\usepackage{xfrac}
\usepackage[unicode,hidelinks,bookmarks=false]{hyperref}
\newcommand{\ie}{i.e.\ }
\newcommand{\cf}{cf.\ }
\newcommand{\resp}{respectively\ }
\newcommand{\Subsection}{\subsection}
\providecommand{\nobreakdash}{\nobreak\hbox{-}}
\setlength{\emergencystretch}{2em}
\multlinegap0pt
\usepackage{bm}
\usepackage{bbm}
\usepackage{dsfont}
\usepackage{xspace}
\usepackage{cancel}
\usepackage{mathtools}
\usepackage{amsthm}
\makeatletter
\@ifundefined{theorem}{\newtheorem{theorem}{Theorem}}{}
\@ifundefined{lemma}{\newtheorem{lemma}[theorem]{Lemma}}{}
\makeatother
\def\om{\omega}
\newcommand{\IR}{\mathbb{R}}
\newcommand{\NN}{\mathbb{N}}
\newcommand{\RR}{\mathbb{R}}
\newcommand{\e}{\varepsilon}
\newcommand{\w}{\omega}
\title[Baire-type properties of topological vector spaces]{Baire-type properties of topological vector spaces}
\author{Saak Gabriyelyan, Alexander V. Osipov, Evgenii Reznichenko}
\date{Source: arXiv:2601.23008v2 (CC BY 4.0)}
\begin{document}
\maketitle

\noindent\emph{Source and licence.} Baire-type properties of topological vector spaces. Version arXiv:2601.23008v2 (CC BY 4.0), distributed under the Creative Commons Attribution 4.0 International licence, as recorded in the arXiv licence metadata for this exact version and shown on the abstract page https://arxiv.org/abs/2601.23008v2. Full rights notice: licenses/arxiv-2601.23008-CC-BY-4.0.txt.\par\smallskip

\noindent\textit{Editorial scope.} Counted unit: the Lemma whose source label is \texttt{l:Rn+1-ls} in the article (source0.tex lines 802--809, source label \texttt{l:Rn+1-ls}), delivered together with the article's own complete original proof of it (source0.tex lines 811--895, one \texttt{proof} environment of 85 lines). No mathematics is written, repaired or re-derived: statement and proof are the article's own text line for line. Required context retained uncounted: none. Every definition, display and label of those lines is retained unchanged. Omitted from this package and not redistributed: 1--801, 810--810, 896--1379. Nothing inside a retained excerpt is omitted or altered: every retained body is delivered exactly as the article's lines are written, so no omission receipt of any kind accompanies this package. Citations: the retained excerpts carry no citation command at all, so no bibliography entry of the article is transcribed and no reference is redistributed. Reference adaptation: none was needed and none was made; every reference inside the delivered unit resolves inside it, so no unresolved reference or citation is produced. Printed slips kept literal: no printed word, symbol, hypothesis or number of the counted statement or of its proof was altered. Layout and class: the article's own class is \texttt{elsarticle} with packages including amsfonts, amsmath, amscd, amssymb, pb-diagram, xy, amsthm, color; the wrapper is the lane's pinned \texttt{amsart} editorial wrapper at 10pt on A4, which restates the theorem-like environments the retained text needs and adds the article's own macro definitions that the retained text needs (\texttt{\textbackslash IR}, \texttt{\textbackslash NN}, \texttt{\textbackslash RR}, \texttt{\textbackslash e}, \texttt{\textbackslash om}, \texttt{\textbackslash w}), with extra wrapper packages (bm, bbm, dsfont, xspace, cancel, mathtools). The delivered numbering is the unit's own. Assets: no retained span contains a figure environment, a graphics-inclusion command, a PDF-page inclusion, a table or a TikZ picture; no image file is copied into this package, the package declares no asset dependency and no image file is redistributed with it. Licence: version arXiv:2601.23008v2 of this article is distributed under the Creative Commons Attribution 4.0 International licence, as recorded in the arXiv licence metadata for this exact version and shown on the abstract page https://arxiv.org/abs/2601.23008v2. A full rights notice is delivered at licenses/arxiv-2601.23008-CC-BY-4.0.txt. Characters: every retained excerpt is pure ASCII, so the delivered engine, XeLaTeX with its default Latin Modern text font, reports no missing character.
\begin{lemma}\label{l:Rn+1-ls}
For every $n\in\NN$, there exists a sequence $(V_m)_{m\in \w}$ of open subsets of $\IR^{n+1}$ such that
\begin{enumerate}
\item[{\rm(i)}] for every $m\in\w$, $V_m$ is bounded and $0\not\in \overline{V_m}$;
\item[{\rm(ii)}] $(V_m)_{m\in \w}$ converges to zero, i.e., for every neighborhood $W$ of $0$ there is $k\in\w$ such that $V_m\subseteq W$ for every $m\geq k$;
\item[{\rm(iii)}] $|\{m\in\w: V_m\cap L\neq\emptyset\}|\leq n $ for any linear subspace $L\subseteq \IR^{n+1}$ of dimension $n$.
\end{enumerate}
\end{lemma}

\begin{proof}
Denote by $\IR_n[x]$ the linear space of all polynomials of degree at most $n$ over the field $\IR$.
We shall identify $\IR^{n+1}$ with the dual space $\IR_n'[x]$ of $\IR_n[x]$. With this identification and taking into account that any $n$-dimensional linear subspace of $\IR_n'[x]$ is uniquely defined by a nonzero polynomial $f\in \IR_n[x]$ (as the kernel of $f$ under the duality $f(l)=l(f)$ for $f\in \IR_n[x]$ and $l\in \IR_n'[x]$), to prove the lemma we first construct a sequence $(U_m)_{m\in \w}$ of open {\em bounded} subsets of $\IR_n'[x]$ such that $0\not\in \overline{U_m}$ for every $m\in\w$, and
\begin{equation}\label{eq:lemRn+1}
|\{m\in\w: 0\in f(U_m) \}|\leq n \; \text{ for each }\; f\in \IR_n[x]\setminus \{0\}.
\end{equation}

For $t\in\IR$, let $\delta_t\in \IR_n'[x]$ denote the Dirac functional: $\delta_t(f)=f(t)$ for $f\in \RR_n[x]$.
\smallskip

{\em {\bf Claim 1.} If $t_0,t_1,...,t_m$ are distinct elements of $\IR$ and $m\leq n$, then the functionals $\delta_{t_0},\delta_{t_1},...,\delta_{t_m}$ are linearly independent.}
\smallskip

{\em Proof of Claim 1.} By adding distinct elements $t_{m+1},\dots,t_n$, without loss of generality we assume that $m=n$. For every $i=0,1,...,n$, let $f_i\in \IR_n[x]$ be such that $f_i(t_j)=0$ if $j\neq i$ and $f_i(t_j)=1$ if $j= i$. Then the matrix $\big(\delta_{t_i}(f_j)\big)_{ij}=\big(f_j(t_i)\big)_{ij}$ is the identity matrix and, therefore, nonsingular. Consequently, the functionals $\delta_{t_0},\delta_{t_1},...,\delta_{t_m}$ are linearly independent. Claim 1 is proved.
\smallskip

We set
\[
\e := \tfrac{1}{2(n+1)2^n} \; \text{ and } \; q_m := (n+1) 2^m \; \text{ for }\; m\in\w.
\]
\smallskip

{\em {\bf Claim 2.} Let $m_0,m_1,...,m_n\in\w$ be distinct numbers, and let $f\in \IR_n[x]$ be such that $|f(q_{m_i})|<\e$ for $i=0,1,...,n$.  Then $|f(k)|\leq \tfrac{1}{2}$ for every $k=0,1,...,n$.}
\smallskip

{\em Proof of Claim 2.} To prove the claim, for every $i=0,1,...,n$, we set $y_i:=f(q_{m_i})$ and
\[
f_i(x) =\frac
{(q_{m_0}-x)(q_{m_1}-x)\cdots(q_{m_{i-1}}-x)(q_{m_{i+1}}-x)\cdots(q_{m_{n}}-x)}
{(q_{m_0}-q_{m_i})(q_{m_1}-q_{m_i})\cdots(q_{m_{i-1}}-q_{m_i})(q_{m_{i+1}}-q_{m_i})\cdots(q_{m_{n}}-q_{m_i})}.
\]
%
Then
\[
f(x)=\sum_{i=0}^n y_i f_i(x).
\]
By the assumption of the claim we have $|y_i|<\e$. Since $n < q_{m_j}$ for each $j=0,1,...,n$, we obtain
\[
|f_i(k)|\leq |f_i(0)| \;\; \mbox{ for every $k=0,1,...,n$}.
\]


Let us estimate $|f_i(0)|$. Since
\[
f_i(0) =\frac{1}
{(1-\frac{q_{m_i}}{q_{m_0}})(1-\frac{q_{m_i}}{q_{m_1}})\cdots(1-\frac{q_{m_i}}{ q_{m_{i-1}}})(1-\frac{q_{m_i}}{q_{m_{i+1}}})\cdots(1-\frac{q_{m_i}}{q_{m_{n}}})}
\]
and $|1-\tfrac{q_{m_i}}{q_{m_j}}|=|1-2^{m_i-m_j}|\geq \frac{1}{2}$ for distinct $i$ and $j$, we obtain
\[
|f_i(0)|\leq 2^{n}.
\]
Therefore, for every $k=0,1,...,n$, we have
\[
|f(k)| = \Big|\, \sum_{i=0}^n y_i f_i(k) \,\Big| \leq \sum_{i=0}^n |y_i| \cdot|f_i(k)| \leq \sum_{i=0}^n |y_i|\cdot |f_i(0)| < \e \cdot (n+1) 2^n
= \tfrac{1}{2}.
\]
Claim 2 is proved.
\smallskip

Now we define
\begin{align*}
V  &:= \big\{ f\in \IR_n[x] : |f(\delta_k)|<1 \text{ for }k=0,1,...,n\big\},
\\
U &:= \big\{ l\in \IR_n'[x] : l(V) \subseteq (-1,1) \big\}.
\end{align*}
Since, by Claim 1,  $\delta_0,\dots,\delta_n$ are linearly independent in the $(n+1)$-dimensional space $\IR'_n[x]$, it follows that $V$ and hence also $U$ are open bounded neighborhoods of zero. For every $m\in\w$, set
\[
U_m := \delta_{q_m} + \e U.
\]
Then all $U_m$ are open bounded subsets of $\IR_n'[x]$. We claim that $0\not\in\overline{U_m}$. Indeed, choose a polynomial $g\in V$ such that $g(q_m)>2$. Then, for every $l\in U$, we have
\[
g(\delta_{q_m} + \e l)\geq g(q_m)-\e>1,
\]
which implies that $0\not\in\overline{U_m}$.

Let us check that the sequence $(U_m)_{m\in\om}$ satisfies condition (\ref{eq:lemRn+1}).
Assuming the converse we could find distinct numbers $m_0,m_1,...,m_n\in\w$ and $g\in \IR_n[x]\setminus\{0\}$ such that $0\in g(U_{m_i})$ for $i=0,1,...,n$.
Let $s=\max \{|g(k)|:k=0,1,...,n\}$. Since $g\neq 0$, we have $s>0$. Set $f(x):= \tfrac{2g(x)}{3s}$.
Then $f\in V$ and $|f(k_0)|>\tfrac{1}{2}$ for some $k_0\in \{0,1,...,n\}$.

Let $i\in\{0,1,...,n\}$. Then $0\in f(U_{m_i})$. Choose $h_i\in U$ such that $f(\delta_{q_{m_i}}-\e h_i)=0$. Then $f(q_{m_i})=\e f(h_i)$. Since $f\in V$ and $h_i\in U$, we have $|f(h_i)|< 1$. Therefore, $|f(q_{m_i})|<\e$. Applying Claim 2 to the polynomial $f$, we obtain  that $|f(k_0)|\leq \tfrac{1}{2}$. This is a contradiction.
\smallskip

To finish the proof of the lemma, for every $m\in\w$,  it suffices to set $V_m:= a_m U_m = a_m(\delta_{q_m} + \e U)$, where $a_m\to 0$ sufficiently fast. Then the properties of $U_m$-s proved above imply (i)-(iii) for the sequence  $(V_m)_{m\in \w}$, as desired.
\end{proof}

\end{document}
